\subsection{THE CONCEPT OF A DIFFERENTIAL EQUATION}

Let I be an interval contained in R, not reducing to a single point, E a topological vector space over R, and A and B two open subsets of E. Let \((\mathbf{x}, \mathbf{y}, t) \mapsto \mathbf{g}(\mathbf{x}, \mathbf{y}, t)\) be a continuous map of \(A \times B \times I\) into E; to every differentiable map u of I into A whose derivative takes its values in B we associate the map \(t \mapsto \mathbf{g}(\mathbf{u}(t), \mathbf{u}'(t), t)\) of I into E, and denote it by \(\tilde{\mathbf{g}}(\mathbf{u})\) ; so \(\tilde{g}\) is defined on the set \(\mathcal{D}(\mathrm{A}, \mathrm{B})\) of differentiable functions of I into B whose derivatives have their values in B. We shall say that the equation \(\tilde{\mathbf{g}}(\mathbf{u}) = 0\) is a differential equation in u (relative to the real variable t); a solution of this equation is also called an integral of the differential equation (on the interval I); it is a differentiable map of I into A, whose derivative takes values in B, such that \(\mathbf{g}(\mathbf{u}(t), \mathbf{u}'(t), t) = 0\) for every \(t \in I\) . By abuse of language we shall write the differential equation \(\tilde{\mathbf{g}}(\mathbf{u}) = 0\) in the form

\[
\mathbf {g} (\mathbf {x}, \mathbf {x} ^ {\prime}, t) = 0,
\]

on the understanding that x belongs to the set D(A, B).

For example, for I = E = R the relations

\[
x ^ {\prime} = 2 t, \qquad t x ^ {\prime} - 2 x = 0, \qquad x ^ {\prime 2} - 4 x = 0, \qquad x - t ^ {2} = 0
\]

are differential equations, all four of which admit the function \(x(t) = t^{2}\) as a solution.

In this chapter we consider in principle only the case where E is a complete normed space over R, and where the differential equations are of the specific form

\[
\mathbf {x} ^ {\prime} = \mathbf {f} (t, \mathbf {x})\tag{1}
\]

(``explicit equations in \(\mathbf{x}'\)''), where \(\mathbf{f}\) denotes a function defined on \(\mathrm{I} \times \mathrm{H}\) with values in E, and H is an open nonempty subset of E. We shall, moreover, widen a little the concept of a solution (or integral) of such an equation (on the interval I): we shall say that a function \(\mathbf{u}\) , defined and continuous on I, with values in H, is a solution (or integral) of the equation (1) if there exists a countable subset A of I such that at every point \(t\) of the complement of A in I the function \(\mathbf{u}\) admits a derivative \(\mathbf{u}'(t)\)

such that \(\mathbf{u}^{\prime}(t)=\mathbf{f}(t,\mathbf{u}(t))\) . If u is differentiable and satisfies the relation above for every point \(t\in I\) we shall say that it is a strict solution of the equation (1) on I.

In the particular case of a differential equation of the form \(\mathbf{x}^{\prime}=\mathbf{f}(t)\) , where f is a map of I into E, the solutions in the above sense are the primitives of the function f (II, p.~51), and the strict solutions are the strict primitives.

When E is a product of complete normed spaces \(E_{i}\) ( \(1 \leqslant i \leqslant n\) ), one can write \(\mathbf{x} = (\mathbf{x}_{i})_{1 \leqslant i \leqslant n}\) and \(\mathbf{f} = (\mathbf{f}_{i})_{1 \leqslant i \leqslant n}\) , where \(x_{i}\) is a map from I into \(E_{i}\) and \(f_{i}\) is a map from \(I \times H\) into \(E_{i}\) ; the equation (1) is then equivalent to the system of differential equations

\[
\mathbf {x} _ {i} ^ {\prime} = \mathbf {f} _ {i} (t, \mathbf {x} _ {1}, \mathbf {x} _ {2}, \dots , \mathbf {x} _ {n}) \quad (1 \leqslant i \leqslant n).\tag{2}
\]

The most important case is that where the \(E_{i}\) are equal to R or to C; one then says that (2) is a system of scalar differential equations.

One may reduce the study of relations of the form

\[
\mathbf {x} ^ {(n)} = \mathbf {f} \big (t, \mathbf {x}, \mathbf {x} ^ {\prime}, \dots , \mathbf {x} ^ {(n - 1)} \big)\tag{3}
\]

to that of the system (2), where \(\mathbf{x}\) is an \(n\) times differentiable vector function on I; for on putting \(\mathbf{x}_1 = \mathbf{x}\) , and \(\mathbf{x}_p = \mathbf{x}^{(p - 1)}\) for \(2 \leqslant p \leqslant n\) , the relation (3) is equivalent to the system

\[
\left\{ \begin{array}{l} \mathbf {x} _ {i} ^ {\prime} = \mathbf {x} _ {i + 1} \\ \mathbf {x} _ {n} ^ {\prime} = \mathbf {f} (t, \mathbf {x} _ {1}, \mathbf {x} _ {2}, \ldots , \mathbf {x} _ {n}). \end{array} \right. \quad (1 \leqslant i \leqslant n - 1)\tag{4}
\]

A relation of the form (3) is called a differential equation of order n (explicitly resolved for \(\mathbf{x}^{(n)}\) ); in contrast, equations of the form (1) are called differential equations of first order.

Similarly one may reduce any ``system of differential equations'' of the form

\[
\mathrm{D} ^ {n _ {i}} \mathbf {x} _ {i} = \mathbf {f} _ {i} (t, \mathbf {x} _ {1}, \mathrm{D} \mathbf {x} _ {1}, \dots , \mathrm{D} ^ {n _ {1} - 1} \mathbf {x} _ {1}, \dots , \mathbf {x} _ {p}, \mathrm{D} \mathbf {x} _ {p}, \dots , \mathrm{D} ^ {n _ {p} - 1} \mathbf {x} _ {p})\tag{5}
\]

\((1 \leqslant i \leqslant p)\) to a system of the form (2), where \(\mathbf{x}_i\) is function on I which is \(n_i\) times differentiable on I (for \(1 \leqslant i \leqslant p\) ).

